\documentclass[10pt,leqno]{article}
\usepackage[T1]{fontenc}
\usepackage[utf8]{inputenc}
\usepackage{lmodern}
\usepackage{amsmath,amssymb,amsthm,mathtools,mathrsfs}
\usepackage[normalem]{ulem}
\usepackage{geometry}
\usepackage{enumitem}
\usepackage{tabularx}
\geometry{margin=0.86in}
\setlength{\parindent}{1.45em}
\setlength{\parskip}{0.18em}
\emergencystretch=2em
\setlist{leftmargin=2.2em,itemsep=0.12em,topsep=0.25em}
\newcommand{\C}{\mathbb C}
\newcommand{\R}{\mathbb R}
\newcommand{\Z}{\mathbb Z}
\newcommand{\Q}{\mathbb Q}
\newcommand{\Pj}{\mathbb P}
\newcommand{\Gm}{\mathbb G_m}
\newcommand{\bS}{\mathbb S}
\newcommand{\cO}{\mathcal O}
\newcommand{\cU}{\mathcal U}
\newcommand{\cC}{\mathcal C}
\newcommand{\Gr}{\operatorname{Gr}}
\newcommand{\Ker}{\operatorname{Ker}}
\newcommand{\Imop}{\operatorname{Im}}
\newcommand{\Res}{\operatorname{Res}}
\newcommand{\QED}{\hfill$\square$}
\newtheoremstyle{deligne}{0.6em}{0.6em}{\itshape}{}{\bfseries}{.-}{0.5em}{\thmname{#1}\thmnote{ #3}}
\theoremstyle{deligne}
\newtheorem*{definition}{Definition}
\newtheorem*{theorem}{Theorem}
\newtheorem*{lemma}{Lemma}
\newtheorem*{proposition}{Proposition}\newtheorem*{corollary}{Corollary}\newtheorem*{monodromyTheorem}{Monodromy Theorem}

\newcommand{\Lie}{\operatorname{Lie}}
\newcommand{\ad}{\operatorname{ad}}
\newcommand{\GL}{\operatorname{GL}}
\newcommand{\Aut}{\operatorname{Aut}}
\newcommand{\End}{\operatorname{End}}
\newcommand{\dd}{\mathrm d}
\renewcommand{\underline}[1]{\uline{#1}}
\renewcommand{\thesection}{\arabic{section}.}

\begin{document}

\begin{center}
Séminaire BOURBAKI\\
22nd year, 1969/70, no. 376\hfill May--June 1970\\[2.2em]
{\Large\bfseries GRIFFITHS'S WORK}\\[1.2em]
by Pierre DELIGNE
\end{center}

This talk contains part of Griffiths's fundamental results on families of Hodge structures. It contains none of his results on algebraic cycles.

\section{\underline{Hodge structures}.}

Let $H_\R$ be a real vector space of finite dimension.

\begin{definition}[1.1]
\underline{A} real Hodge structure \underline{on} $H_\R$ \underline{is a bigrading of}
\[
H_\C=H_\R\otimes_\R\C
\]
\[
H_\C=\bigoplus_{p,q}H^{pq},
\]
\underline{such that} $H^{pq}$ \underline{and} $H^{qp}$ \underline{are complex conjugates}.
\end{definition}

\underline{The Hodge numbers} of a real Hodge structure $H$ are the integers
\[
h^{pq}=\dim_\C(H^{pq}).
\]
One says that $H$ has \underline{weight} $n$ if $h^{pq}=0$ for $p+q\ne n$. A variant of 1.1 is the following.

\begin{definition}[1.2]
\underline{A} real Hodge structure of weight $n$ \underline{on} $H_\R$ \underline{is a finite decreasing filtration} $F$ \underline{of} $H_\C$ (\underline{the} Hodge filtration) \underline{such that}, \underline{for} $p+q=n+1$, \underline{one has}
\begin{equation}
F^p(H_\C)\oplus \overline F^{q}(H_\C)\xrightarrow{\sim}H_\C .
\tag{1.2.1}
\end{equation}
\end{definition}

One passes from 1.1 to 1.2 by putting
\begin{equation}
F^p(H_\C)=\sum_{p'\ge p}H^{p'q'}
\tag{1.2.2}
\end{equation}
and one passes from 1.2 to 1.1 by putting, for $p+q=n$,
\begin{equation}
H^{pq}=F^p(H_\C)\cap \overline F^{q}(H_\C) .
\tag{1.2.3}
\end{equation}

\paragraph{1.3.}
Let $\bS$ denote the real algebraic group obtained from the multiplicative group $\Gm$ by Weil restriction of scalars from $\C$ to $\R$. We have $\bS(\R)=\C^*$; there is a natural morphism $w$ from $\Gm$ to $\bS$ which, on real points, is the inclusion of $\R^*$ into $\C^*$, and there is a natural morphism $t$ from $\bS$ to $\Gm$ which, on real points, is the norm. The composite $tw$ is the squaring map. A new variant of 1.1 is the following.

\begin{definition}[1.4]
\underline{A real Hodge structure} \underline{on} $H_\R$ \underline{is an action} (\underline{defined over} $\R$) \underline{of the algebraic group} $\bS$ \underline{on} $H_\R$.
\end{definition}

If $z$ denotes the defining isomorphism $z:\bS(\R)\xrightarrow{\sim}\C^*$, one passes from 1.1 to 1.4 by defining $H^{pq}$ as the subspace of $H_\C$ on which $\bS(\R)$ acts by multiplication by $z^p\overline z^{\,q}$.

If $H_\R$ is endowed with a real Hodge structure, the Weil automorphism $C$ is defined as multiplication by $i^{p-q}$ on $H^{pq}$. This automorphism is real and is simply the action of the element $i$ of $\bS(\R)\subset\C^*$.

\begin{definition}[1.5]
\underline{A} \underline{polarization of a real Hodge structure} $H$ \underline{of weight} $n$ \underline{is a bilinear form} $\Psi$ \underline{on} $H_\R$, \underline{invariant under the compact subgroup} $\Ker(t)$ \underline{of} $\bS$, \underline{and such that the form} $\Psi(x,Cy)$ \underline{is symmetric and positive definite}.
\end{definition}

From the identity $C^2=(-1)^n$, one obtains
\[
\Psi(x,y)=(-1)^n\Psi(x,C^2y)=(-1)^n\Psi(Cy,Cx)=(-1)^n\Psi(y,x)
\]
so that $\Psi$ is alternating or symmetric according to the parity of $n$. The variance property under $\bS$ and the positivity property are written as follows:
\begin{equation}
\text{the orthogonal complement of }F^p(H_\C)\text{ with respect to }\Psi\text{ is }F^{n-p}(H_\C);\tag{1.5.1}
\end{equation}
\begin{equation}
\begin{aligned}
&\text{for nonzero }x\text{ in }H^{pq},\text{ one has}\\
&i^{p-q}\Psi(x,\overline x)>0
\end{aligned}
\tag{1.5.2}
\end{equation}
(Riemann bilinear relations).

\begin{definition}[1.6]
(i) \underline{A} Hodge structure $H$ of weight $n$ \underline{consists of a finitely generated free} $\Z$\underline{-module} $H_\Z$ (\underline{the ``integral lattice''}) \underline{and a real Hodge structure of weight} $n$ \underline{on} $H_\R=H_\Z\otimes_\Z\R$.

(ii) \underline{A} \underline{polarization of a Hodge structure} $H$ \underline{of weight} $n$ \underline{is a polarization} $\Psi$ \underline{of the underlying real Hodge structure} \underline{that takes integral values on the integral lattice} $H_\Z$.
\end{definition}

\section{\underline{Hodge theory}.}

\paragraph{2.1.}
Let $X\subset\Pj^r(\C)$ be a nonsingular projective complex algebraic variety, pure of dimension $d$, and let $\eta\in H^2(X,\Z)$ be the cohomology class of a hyperplane section of $X$, i.e., the first Chern class of $\cO(1)$.

The holomorphic de Rham complex $\Omega_X^*$ is a resolution of the constant sheaf $\C$ (the holomorphic Poincaré lemma), whence a spectral sequence
\begin{equation}
E_1^{pq}=H^q(X,\Omega_X^p)\Rightarrow H^{p+q}(X,\C).
\tag{2.1.1}
\end{equation}

\paragraph{2.2.}
Hodge theory asserts that
\begin{enumerate}[label=(\Alph*)]
\item the spectral sequence (2.1.1) degenerates ($E_1=E_\infty$).
\item the filtration of $H^n(X,\C)$ to which it abuts is a Hodge structure of weight $n$ (1.2) on $H^n(X,\R)$.
\end{enumerate}
By (A), one has
\[
H^{pq}\simeq H^q(\Omega_X^p).
\]

(C) For $n\le d$, let $P_\Q^n$ denote the kernel of the iterated cup-product
\[
\eta^{d-n+1}:H^n(X,\Q)\longrightarrow H^{2d-n+2}(X,\Q)
\]
(\underline{the primitive part of cohomology}), and let $P_\Z^n$ be the intersection of $P_\Q^n$ with the image of $H^n(X,\Z)$ in $H^n(X,\Q)$. Since the class $\eta$ is of type $(1,1)$, $P_\R^n=P_\Z^n\otimes\R$ is a Hodge substructure of $H^n(X,\R)$. Up to a sign depending only on $d$ and $n$, the form
\[
\Psi(x,y)=\int \eta^{d-n}\wedge x\wedge y
\]
is a polarization of the Hodge structure $P^n$.

(D) The natural map
\[
\bigoplus_{n-2k\le d}\eta^k\wedge:
\bigoplus P_\Q^{\,n-2k}\longrightarrow H^n(X,\Q)
\]
is an isomorphism (\underline{the Hodge-Lepage decomposition}).

\section{\underline{Families of Hodge structures}.}

\paragraph{3.1.}
Let $f:X\to S$ be a projective and smooth morphism of analytic spaces, of pure relative dimension $d$. For simplicity, we suppose that $S$ is nonsingular and that $f$ is provided with a factorization through $\Pj^r_S=S\times\Pj^r(\C)$. The fibers $X_s=f^{-1}(s)$ therefore form an analytic family, parametrized by $S$, of nonsingular subvarieties of $\Pj^r(\C)$.

\paragraph{3.2.}
From the $C^\infty$ point of view, $f$ is a locally trivial fibration. The cohomology algebra of the fibers of $f$ therefore forms a local system on $S$, denoted
\[
R_\Z^*(f)=\sum R^n f_*\Z .
\]
For every abelian sheaf $F$ on $S$ (for example $\Q$, $\C$, $\cO$), put
\begin{equation}
R_F^n(f)=F\otimes R_\Z^n(f).
\tag{3.2.1}
\end{equation}
Since $f$ is proper, one also has
\[
R_F^n(f)\xrightarrow{\sim} R^n f_*(f^*F).
\]

The cohomology classes of the hyperplane sections of the $X_s$ define a \underline{locally constant} section
\begin{equation}
\eta\in H^0(S,R_\Z^2(f)).
\tag{3.2.2}
\end{equation}
For $n\le d$, the $P_\Z^n(X_s)$ therefore form a local subsystem $P_\Z^n(f)$ of $R_\Q^n(f)$. Putting $P_F^n(f)=F\otimes P_\Z^n(f)$, one has by definition
\begin{equation}
P_\Q^n(f)=\Ker\bigl(\eta^{d-n+1}\wedge:R_\Q^n(f)\to R_\Q^{2d-n+2}(f)\bigr),
\tag{3.2.3}
\end{equation}
\begin{equation}
P_\Z^n(f)=P_\Q^n(f)\cap \Imop\bigl(R_\Z^n(f)\to R_\Q^n(f)\bigr).
\tag{3.2.4}
\end{equation}
The polarization form of 2.2 (C)
\begin{equation}
\Psi(x,y)=\pm\int_{X_s}\eta^{d-n}\wedge x\wedge y
\tag{3.2.5}
\end{equation}
is locally constant on $S$, i.e., it defines
\begin{equation}
\Psi:P_\Z^n(f)\otimes P_\Z^n(f)\longrightarrow \Z .
\tag{3.2.6}
\end{equation}

\paragraph{3.3.}
For $F$ a coherent analytic sheaf on $S$, $f^\cdot F$ will denote its inverse image as a sheaf, and
\[
f^*F=\cO_X\otimes_{f^\cdot\cO_S}f^\cdot F
\]
its inverse image as a sheaf of modules. For $s\in S$, $F_{(s)}$ denotes the $\cO_{S,s}$-module of germs of sections of $F$ at $s$; the fiber of $F$ at $s$ is the vector space
\[
F_s=F_{(s)}\otimes_{\cO_{S,s}}\C .
\]
Recall finally that the relative de Rham complex $\Omega_{X/S}^*$ is the quotient of $\Omega_X^*$ with components $\Omega^1_{X/S}=\Omega_X^1/f^*\Omega_S^1$ and $\Omega^p_{X/S}=\bigwedge^p\Omega_{X/S}^1$.
If $T^1_{X/S}$, the dual of $\Omega^1_{X/S}$, is the relative tangent bundle, then one has ``contraction'' pairings
\begin{equation}
\lrcorner:T^1_{X/S}\otimes\Omega^p_{X/S}\longrightarrow\Omega^{p-1}_{X/S}.
\tag{3.3.1}
\end{equation}

\paragraph{3.4.}
By the relative holomorphic Poincaré lemma, $\Omega_{X/S}^*$ is a resolution of $f^\cdot\cO_S$. Hence one has a spectral sequence
\begin{equation}
E_1^{pq}=R^qf_*\Omega^p_{X/S}\Rightarrow R^n f_*(f^\cdot\cO_S)=R_\cO^n(f).
\tag{3.4.1}
\end{equation}
It follows from 2.2 (A) that the $E_1^{pq}$ are locally free, that the spectral sequence (3.4.1) degenerates ($E_1=E_\infty$) and that its formation commutes with arbitrary base change. The filtration of $R_\cO^n(f)$ to which it abuts (\underline{the Hodge filtration}) therefore induces on $R_\cO^n(f)_s\simeq H^n(X_s,\C)$ the Hodge filtration of 2.2: the latter varies holomorphically with $s\in S$.

\begin{definition}[3.5]
\underline{The Gauss-Manin connection on the holomorphic vector bundle} $R_\cO^n(f)$ (\underline{identified here with its sheaf of holomorphic sections}) \underline{is the holomorphic and integrable connection}
\[
\nabla:R_\cO^n(f)\longrightarrow\Omega_S^1\otimes R_\cO^n(f)
\]
\underline{whose horizontal local sections} ($\nabla v=0$) \underline{are the local sections of} $R_\C^n(f)$.
\end{definition}

\begin{theorem}[of transversality 3.6]
\underline{The Hodge filtration} $F$ \underline{on} $R_\cO^n(f)$ \underline{satisfies}
\[
\nabla F^p(R_\cO^n(f))\subset \Omega_S^1\otimes F^{p-1}(R_\cO^n(f)).
\]
\end{theorem}

\underline{Proof} (following Katz-Oda [10] and a personal communication from Katz). The question is local on $S$. Let $\cU=(U_i)_{0\le i\le N}$ be a finite covering of $X$ by Stein opens. For $Q\subset[0,N]$, let $U_Q=\bigcap_{i\in Q}U_i$, let $j_Q$ be the inclusion of $U_Q$ in $X$, and put
\[
f_*(U_Q,\Omega^p_{X/S})=(fj_Q)_*(\Omega^p_{X/S}|U_Q).
\]
We denote by $f_*(\cU,\Omega^*_{X/S})$ the double complex of sheaves on $S$ with components
\begin{equation}
f_*(\cU,\Omega^*_{X/S})^{pq}=\bigoplus_{\#Q=q+1} f_*(U_Q,\Omega^p_{X/S})
\tag{3.6.1}
\end{equation}
whose first differential is induced by the exterior differential and whose second differential is the ``Čech'' differential. The spectral sequence 3.4.1 is the spectral sequence of this double complex obtained from the filtration by the first degree.

Let $v$ be a (holomorphic) vector field on $S$, and suppose that on each $U_i$ a vector field $v_i$ lifting $v$ is given. We denote by $\theta(v_i)$ the endomorphism of $f_*(\cU,\Omega^*_{X/S})$ satisfying, for $Q=(i_1,\ldots,i_q)$ with $i_1<\cdots<i_q$,
\[
\theta(v_i) f_*(U_Q,\Omega^p_{X/S})
\subset
f_*(U_Q,\Omega^p_{X/S})
\oplus
\bigoplus_{i_0<i_1} f_*(U_{\{i_0\}\cup Q},\Omega^{p-1}_{X/S}),
\]
and having as coordinates the Lie derivative
\[
\mathcal L_{v_{i_1}}:f_*(U_Q,\Omega^p_{X/S})\longrightarrow f_*(U_Q,\Omega^p_{X/S})
\]
and the contractions
\[
(-1)^p(v_{i_1}-v_{i_0})\lrcorner:
 f_*(U_Q,\Omega^p_{X/S})\longrightarrow f_*(U_{\{i_0\}\cup Q},\Omega^{p-1}_{X/S}).
\]
A simple calculation shows that $\theta(v_i)$ commutes with $d$.

\begin{lemma}[3.7]
\underline{On the cohomology sheaves} $R_\cO^n(f)$ of $f_*(\cU,\Omega^*_{X/S})$, \underline{the endomorphism} $\theta(v_i)$ \underline{induces} $\nabla_v$.
\end{lemma}

Let $\cC^\infty$ be the sheaf of $C^\infty$ functions on $S$. If $f_*(\cU,\Omega^*_{\infty X/S})$ is the $C^\infty$ analogue of (3.6.1), constructed in terms of the complex of relative $C^\infty$ differential forms with complex values on $X$, then
\[
R_{\cC^\infty}^n(f)=\mathcal H^n(f_*(\cU,\Omega^*_{\infty X/S})).
\]
Since the sheaves $\Omega^p_{\infty X/S}$ are soft, this formula remains valid for every covering $\cU'$.

If $(v_i')_{0\le i\le N}$ are $C^\infty$ liftings of $v$ on the $U_i$, the construction 3.6 still defines endomorphisms $\theta(v_i')$ of $R_{\cC^\infty}^n(f)$. If $(v_i')$ and $(v_i'')$ are two systems of liftings, a simple calculation shows that, at the level of complexes,
\begin{equation}
\theta(v_i')-\theta(v_i'')=dH-Hd,
\tag{3.7.1}
\end{equation}
where the only nonzero coordinates of $H$ are the contractions
\[
v'_{i_1}-v''_{i_1}\lrcorner:
 f_*(U_Q,\Omega^p_{\infty X/S})\longrightarrow f_*(U_Q,\Omega^{p-1}_{\infty X/S}),
\]
$(Q=\{i_1,\ldots,i_q\},\ i_1<\cdots<i_q)$.

Thus the endomorphism $\theta(v_i')$ of $R_{\cC^\infty}^n(f)$ is independent of the choice of the $v_i'$. Via the injection of $R_\cO^n(f)$ into $R_{\cC^\infty}^n(f)$, it is compatible with the endomorphism $\theta(v_i)$ of 3.7. It therefore remains only to verify that $\theta(v_i')=\nabla_v$ \underline{on} $R_{\cC^\infty}^n(f)$. One verifies that $\theta(v_i')$, which is independent of the $v_i'$, is also independent of $\cU$. Taking $\cU=\{X\}$ and a lift $v'$ of $v$, one then has
\[
f_*(\cU,\Omega^*_{\infty X/S})=f_*(\Omega^*_{\infty X/S}),
\]
and $\theta(v')$ is the Lie derivative along $v'$, visibly equal to $\nabla_v$.

\paragraph{3.8.}
We finish the proof of 3.6. By construction, one has
\begin{equation}
\theta(v_i)\left(\bigoplus_{p'\ge p} f_*(\cU,\Omega^*_{X/S})^{p',q'}\right)
\subset
\bigoplus_{p'\ge p-1} f_*(\cU,\Omega^*_{X/S})^{p',q'} .
\tag{3.8.1}
\end{equation}
Since the spectral sequence (3.4.1) is obtained from the double complex $f_*(\cU,\Omega^*_{X/S})$, it follows from 3.7 and (3.8.1) that
\begin{equation}
\theta(v_i)F^p(R_\cO^n(f))=\nabla_vF^p(R_\cO^n(f))\subset F^{p-1}(R_\cO^n(f)).
\tag{3.8.2}
\end{equation}
Since this is true locally on $S$ and for every $v$, it implies 3.6.

\paragraph{3.9.}
By the Leibniz formula
\[
\nabla(fh)=df\cdot h+f\cdot \nabla h
\]
the map induced by $\nabla$
\[
\operatorname{def}(\nabla):\Gr_F^p(R_\cO^n(f))\longrightarrow
\Omega_S^1\otimes\Gr_F^{p-1}(R_\cO^n(f))
\]
is $\cO$-linear. By 3.4.1 it is identified with
\begin{equation}
\operatorname{def}(\nabla):R^qf_*\Omega^p_{X/S}\longrightarrow
\Omega_S^1\otimes R^{q+1}f_*\Omega^{p-1}_{X/S}.
\tag{3.9.1}
\end{equation}
It follows from formula 3.7 for $\nabla_v$ that

\begin{proposition}[3.10]\leavevmode\par\noindent
\underline{The homomorphism} (3.9.1) \underline{is the cup product}, \underline{via the pairing} (3.3.1), \underline{with the Kodaira--Spencer class}
\[
c\in\Omega_S^1\otimes R^1f_*(T^1_{X/S})
\]
\underline{which expresses how} $X_s$ \underline{deforms with} $s$.
\end{proposition}

\begin{definition}[3.11]
(i) \underline{A} family $H$ of real Hodge structures \underline{of weight} $n$ \underline{on} $S$ \underline{consists of}
\begin{enumerate}[label=\alph*)]
\item \underline{a local system of real vector spaces} $H_\R$ \underline{on} $S$;
\item \underline{a finite holomorphic filtration} $F$ \underline{of the locally free analytic sheaf}
\[
H_\cO=H_\R\otimes_\R\cO,
\]
\end{enumerate}
\underline{these data satisfying the following conditions}:
\begin{enumerate}[label=(H.\arabic*)]
\item \underline{The natural connection} $\nabla$ \underline{on} $H_\cO$ \underline{is such that}
\[
\nabla F^p(H_\cO)\subset \Omega_S^1\otimes F^{p-1}(H_\cO);
\]
\item \underline{At every point} $s\in S$, $F$ \underline{defines on} $(H_\R)_s$ \underline{a Hodge structure of weight} $n$.
\end{enumerate}
(ii) \underline{A} polarization \underline{of} $H$ \underline{is a locally constant bilinear form} $\Psi$ \underline{on} $H_\R$ \underline{which at every point} $s\in S$ \underline{induces a polarization of} $(H_\R)_s$.
\end{definition}

\paragraph{3.12.}
\underline{A family of Hodge structures} $H$ of weight $n$ on $S$ is a local system $H_\Z$ of finitely generated free $\Z$-modules on $S$, together with a family of real Hodge structures on $H_\R=\R\otimes H_\Z$. A \underline{polarization} of $H$ is a polarization $\Psi$ of $H_\R$ taking integral values on $H_\Z$.

With these definitions, one can reformulate 2.2 (C) and 3.6 by saying that, under the hypotheses of 3.1, $P_\Z^n(f)$ \underline{is a polarized family of Hodge structures} on $S$.

\section{\underline{The regularity theorem}.}

\paragraph{4.1.}
When one starts with a projective and smooth morphism $f:X\to S$ of \underline{schemes} of finite type over $\C$, such that $f^{\mathrm{an}}:X^{\mathrm{an}}\to S^{\mathrm{an}}$ satisfies the hypotheses of 3.1, some of the objects constructed in no. 3 are obtained from purely algebraic objects by applying the ``passage to the analytic category'' functor. They are:
\begin{enumerate}[label=\alph*)]
\item the sheaves of modules $R_\cO^n(f)$, their Hodge filtration, the spectral sequence 3.4.1, and the Gauss--Manin connection;
\item the algebra structure (cup-product) on $R_\cO^*(f)$, and the ``primitive part'' $P_\cO^n(f)$;
\item the product of the polarization form $\Psi$ by a suitable power of $2\pi i$.
\end{enumerate}
By contrast, the ``integral lattice'' $R_\Z^*(f)$, and the real structure thereby obtained on the complex cohomology of the fibers of $f$, are transcendental in nature.

\paragraph{4.2.}
Indeed, let $S$ be a smooth scheme of finite type over $\C$, and let $X$ be a closed subscheme of $\Pj^r(\C)\times S$, such that the projection $f:X\to S$ is a smooth morphism of pure relative dimension $d$.

By the relative variant of GAGA, for every coherent algebraic sheaf $F$ on $X$ one has
\[
(R^nf_*F)^{\mathrm{an}}\xrightarrow{\sim}R^nf_*^{\mathrm{an}}(F^{\mathrm{an}}).
\]
More generally, if $K^*$ is a complex of coherent algebraic sheaves whose


differentials $d_i$ are (algebraic) differential operators which are $f^\cdot\cO_S$-linear, then the direct images in hypercohomology satisfy
\[
  \bigl(R^n f_*(K^*)\bigr)^{\mathrm{an}}
  \xrightarrow{\sim}
  R^n f_*^{\mathrm{an}}(K^{*\mathrm{an}})
\]
(this case reduces to the preceding one \underline{via} one of the hypercohomology spectral sequences).

Take $K$ to be the relative algebraic de Rham complex $\Omega^*_{X/S}$. One obtains (in the notation of 3.2)
\[
\bigl(R^n f_*\Omega^*_{X/S}\bigr)^{\mathrm{an}}
 \xrightarrow{\sim}
 R^n f_*^{\mathrm{an}}\bigl((\Omega^*_{X/S})^{\mathrm{an}}\bigr)
 \xleftarrow{\sim}
 R^n f_*^{\mathrm{an}}\bigl(f^{\cdot\mathrm{an}}\cO_S^{\mathrm{an}}\bigr)
 =
 R^n_{\cO}(f).
\]
Moreover, the spectral sequence (3.4.1) comes from a purely algebraic spectral sequence
\[
\tag{4.2.1}
E_1^{pq}=R^q f_*\Omega^p_{X/S}
\Longrightarrow
R^{p+q}f_*\Omega^*_{X/S}.
\]
According to 3.4, the $E_1^{pq}$ are locally free (since the $E_1^{pq\,\mathrm{an}}$ are), the spectral sequence (4.2.1) degenerates, and its formation is compatible with every change of base. The Hodge filtration to which it abuts is therefore, \underline{ipso facto}, algebraic.

The cup-product is defined in terms of the exterior product in $\Omega^*_{X/S}$. The de Rham cohomology class of a hyperplane section can be defined algebraically (up to multiplication by a power of $2\pi i$), from the logarithmic differential
\[
  \dd f/f:\cO_X^*\longrightarrow \Omega^1_{X/S},
\]
the primitive part $P^n_{\cO}(f)$ and $\Psi$ are also algebraic in nature.

Finally, the construction 3.6--3.7 of the Gauss--Manin connection carries over easily to the algebraic case (Katz--Oda [10]).

\paragraph{4.3.}
Let $D$ be a Riemann surface isomorphic to the unit disk, let $0$ be a point of $D$, put $D^*=D-\{0\}$, let $V$ be a holomorphic vector bundle on $D^*$, let $\nabla$ be a holomorphic connection on $V$, and let $V_o$ be the local system of horizontal sections of $V$. One has
\[
  V \xrightarrow{\sim} V_o\otimes_\C \cO .
\]
The (commutative) fundamental group $\pi_1(D^*)\simeq \Z$ acts on $V_o$. We denote \underline{by} $T$ the monodromy transformation, i.e. the action on $V_o$ or on $V$ of the \underline{positive} generator of $\pi_1(D^*)$.

There exists $U$ such that
\[
  T=\exp(2\pi iU).
\]
Let $z:D\to\C$ be a coordinate, with $z(0)=0$, and let $\widetilde D^*$ be the (universal) covering of $D^*$ on which $\log z$ is defined. For $v$ a section of $V_o$ on $\widetilde D^*$,
\[
  \exp(-\log z\cdot U)(v)
\]
is the inverse image of a holomorphic section of $V$ on $D^*$.

This construction defines an isomorphism
\[
  m_{U,z}:\cO\otimes_\C H^0(\widetilde D^*,V_o)
  \xrightarrow{\sim} V,
\]
and hence a natural extension $V_U$ of $V$ to $D$, depending only on $U$. A section $v$ of $V$ on $D^*$ will be called $\nabla$-moderate if, as a section of $V_U$, it is meromorphic at $0$. This notion does not depend on the choice of $U$.

\paragraph{4.4.}
Let $S$ be the complement, in a smooth projective curve $\overline S$, of a finite set $T$, and let $V$ be an algebraic vector bundle on $S$. An (algebraic) connection $\nabla$ \underline{on} $V$ is called regular if, for every $t\in T$, the meromorphic sections of $V$ in a neighborhood of $t$ are $\nabla$-moderate.

\begin{theorem}[4.5]
\underline{For} $f:X\to S$ \underline{a projective and smooth morphism of schemes}, \underline{the Gauss--Manin connection} on $R^n f_*\Omega^*_{X/S}$ \underline{is regular}.
\end{theorem}

This theorem is proved in [2]. A completely different proof, proceeding by arithmetic methods, is due to Katz [9].

\paragraph{4.6.}
Let $D^*$ be as in 4.3, let $H$ be a family of real Hodge structures on $D^*$, let $T$ be the monodromy transformation of $H_\C$, let $U$ be an endomorphism of $H_\C$ such that $T=\exp(2\pi iU)$, and let $H_{\cO,U}$ be the corresponding locally free extension of $H_\cO$ to $D$.

One says that $H$ is \underline{regular} at $0$ if the Hodge filtration of $H_\cO$ extends to $H_{\cO,U}$. This condition does not depend on $U$. Theorem 4.5 has the following corollary, whose conclusion is purely analytic.

\begin{corollary}[4.7]
\underline{Under the hypotheses of 4.5}, \underline{the families of Hodge structures} $R^n(f)$ \underline{on} $S^{\mathrm{an}}$ \underline{are regular in a neighborhood of every point} $t\in T$.
\end{corollary}

One may hope that, in fact, every polarized family of Hodge structures on $D^*$ is automatically regular.

\section{\underline{Moduli Spaces}.}

\paragraph{5.1.}
Let $G$ be a real algebraic group, let $C\in G(\R)$, and let $V$ be a real representation of $G$. A $C$\underline{-polarization} of $V$ is a bilinear form $\Psi$ on $V$, invariant under $G$, and such that $\Psi(x,Cy)$ is symmetric and positive definite.

\begin{lemma}[5.2]
\underline{If} $G$ \underline{is connected} (\underline{by which one means that} $G(\C)$ \underline{is connected}), \underline{the following conditions are equivalent}:
\begin{enumerate}[label=(\roman*)]
\item $G$ \underline{admits a representation} $\rho:G\to \GL(V)$, \underline{with finite kernel} $\Ker(\rho)$, \underline{which is} $C$\underline{-polarizable}.
\item \underline{Every representation of} $G$ \underline{is} $C$\underline{-polarizable}.
\item $G$ \underline{is reductive} and $\ad C$ \underline{is a Cartan involution} (\underline{of} $G$).
\end{enumerate}
\end{lemma}

The proof is left to the reader.

\paragraph{\underline{Remark} 5.3.-}
There exist in $G(\R)$ elements $C$ satisfying the conditions of 5.2 if and only if $G$ is reductive and admits a compact maximal torus. In this case:
\begin{enumerate}[label=\alph*)]
\item Such an element $C$ is contained in a single maximal compact subgroup, namely its centralizer.
\item For $G$ adjoint, the correspondence $C\longmapsto (\text{centralizer of }C)$ is a bijection between the set of these elements and the set of maximal compact subgroups of $G(\R)$.
\end{enumerate}

\paragraph{5.4.}
Let $G$ be a real algebraic group, endowed with homomorphisms
\[
\tag{5.4.1}
 \mathbb G_m \xrightarrow{w} G \xrightarrow{t} \mathbb G_m
\]
such that $w$ is central and $tw(x)=x^2$. A representation $\rho:G\to\GL(V)$ will be called \underline{homogeneous of weight} $n$ if $\rho w(x)=x^n$.

\underline{By a morphism from} $\bS$ to $G$ we shall always mean a homomorphism $h:\bS\to G$ making the diagram
\[
\begin{array}{ccccc}
\mathbb G_m&\xrightarrow{w}&\bS&\xrightarrow{t}&\mathbb G_m\\
\Vert&&\downarrow h&&\Vert\\
\mathbb G_m&\xrightarrow{w}&G&\xrightarrow{t}&\mathbb G_m
\end{array}
\]
commutative.

If $h:\bS\to G$ is a morphism and $\rho:G\to\GL(V)$ a real representation of $G$, we shall denote by $h_V$ the real Hodge structure (1.4) $\rho\circ h$ on $V$. For $\rho$ homogeneous of weight $n$, a \underline{polarization} of $V$ is a polarization $\Psi$ of $(V,h_V)$ which satisfies
\[
\tag{5.4.2}
  \Psi(gx,gy)=t(g)^n\Psi(x,y).
\]
We shall say that $h$ is \underline{positive} if every representation of $G$ is polarizable.

\paragraph{5.5.}
The group $G$, endowed with (5.4.1), is uniquely determined by $G^o=\Ker(t)$, endowed with the central element, of order dividing $2$, $\varepsilon=w(-1)$. Indeed $G$ is identified with the quotient of $\mathbb G_m\times G^o$ by the subgroup generated by $(-1,\varepsilon)$.

From this point of view, a morphism $h:\bS\to G$ is identified with $h^o:\bS^o\to G^o$ such that $h^o(-1)=\varepsilon$. A homogeneous representation of weight $n$ of $G$ is identified with a representation $\rho$ of $G^o$ such that $\rho(\varepsilon)=(-1)^n$, and a polarization of this representation is identified with an $h^o(i)$-polarization (5.1). According to 5.2, if $G^o$ is connected, $h$ is a positive morphism if and only if $G$ is reductive and $\ad h^o(i)$ is a Cartan involution of $G^o$.

\paragraph{5.6.}
\underline{We shall henceforth suppose that} $G^o$ \underline{is reductive and connected}.

Let $h$ be a morphism from $\bS$ to $G$. For $\rho:G\to\GL(V)$ a representation of $G$, denote by $F_h$ the Hodge filtration of $V_\C$ deduced from $h_V$. For the adjoint representation $\Lie(G)$,
\[
\tag{5.6.1}
  F_h^0(\Lie(G)_\C)=\bigoplus_{p\ge0}\Lie(G_\C)^{p,-p}
\]
is the Lie algebra of a parabolic subgroup $P(h)$ of $G_\C$. For every representation $\rho:G\to\GL(V)$ of $G$, $P(h)$ respects the Hodge filtration $F_h$ of $V_\C$; if $\Ker(\rho)\cap G^o$ is central, then $P(h)$ is exactly the subgroup of $G$ which respects this filtration.

Finally, denote by $H(h)$ the centralizer of $h$ in $G(\R)$.

\paragraph{5.7.}
Let $X$ be a conjugacy class of morphisms from $\bS$ to $G$. We shall denote by $\check X$ the set of conjugates in $G(\C)$ of $P(h)$, for $h\in X$.

\begin{lemma}[5.8]
\underline{The $G(\R)$-equivariant map} $P:h\mapsto P(h)$ \underline{identifies} $X$ \underline{with an open subset of the flag space} $\check X$.
\end{lemma}

Let $h\in X$ and let $\rho:G\to\GL(V)$ be a homogeneous representation of $G$ such that $\Ker(\rho)\cap G^o$ is finite. The set $X$ is identified with the set of Hodge structures on $V$, conjugate under $G(\R)$ to $h_V$, while $\check X$ is identified with the set of filtrations on $V_\C$, conjugate under $G(\C)$ to $F_h$. Since a homogeneous Hodge structure is determined by its Hodge filtration, the map $P$ is injective. It is identified with the map
\[
  G(\R)/H(h)\longrightarrow G(\C)/P(h)
\]
whose differential at the origin is
\[
  \Lie(G)/\Lie(H)\longrightarrow \Lie(G_\C)/\Lie(P(h)),
\]
i.e.
\[
  \Lie(G)/(F^0\cap\overline{F^0})(\Lie(G))
  \longrightarrow
  F^{-\infty}/F^0(\Lie(G_\C)).
\]
This differential is bijective, whence 5.8.

We have seen in passing that
\[
\tag{5.8.1}
  H(h)=G(\R)\cap P(h).
\]

\paragraph{5.9.}
Let $V$ be a real representation of $G$, let $V(\check X,\R)$ be the constant sheaf on $\check X$ with value $V$, put $V(\check X,\C)=\C\otimes V(\check X,\R)$ and $V(\check X,\cO)=\cO\otimes_\C V(\check X,\C)$, and let $\nabla$ be the natural integrable connection of $V(\check X,\cO)$. The $G(\R)$-equivariant map $h\mapsto h_V$ from $X$ to the Hodge structures on $V$ extends uniquely to a $G(\C)$-equivariant, and hence holomorphic, map from $\check X$ to the filtrations of $V_\C$. This map corresponds to a filtration $F$ of $V(\check X,\cO)$, \underline{the Hodge filtration}. The system $(V(\check X,\cO),\nabla,F)$ is $G(\C)$-equivariant. The real structure $V(\check X,\R)$ of $V(\check X,\C)$ is $G(\R)$-equivariant. Finally, at $h\in X$, $F$ induces on $V_\C$ the Hodge filtration $F_h$.

\paragraph{5.10.}
The action of $G(\C)$ on $\check X$ defines a morphism
\[
\tag{5.10.1}
  \Lie(G)(\check X,\cO)\longrightarrow T_{\check X},
\]
where $T_{\check X}$ is the tangent bundle. At $h\in X$, the kernel of (5.10.1) is $\Lie(P(h))=F_h^0(\Lie(G_\C))$ (5.6), so that (5.10.1) factors through a $G(\C)$-equivariant isomorphism
\[
\tag{5.10.2}
  F^{-\infty}/F^0\bigl(\Lie(G)(\check X,\cO)\bigr)
  \xrightarrow{\sim} T_{\check X}.
\]

Let $T^1_{\check X}$ denote the holomorphic subbundle of $T_{\check X}$ which is the image of $F^{-1}(\Lie(G)(\check X,\cO))$. It is easily verified that if $v$ is a local section of $T^1_{\check X}$, then
\[
  \nabla_v\bigl(F^p(V(\check X,\cO))\bigr)
  \subset F^{p-1}(V(\check X,\cO)),
\]
whatever the representation $\rho:G\to\GL(V)$ may be. The converse is true when $\Ker(\rho)\cap G^o$ is central.

\paragraph{5.11.}
Suppose now that $X$ is a conjugacy class of \underline{positive} morphisms from $\bS$ to $G$. For $h\in X$, put $C_h=h(i)$, and denote by $K(h)$ the centralizer of $C_h$ in $G(\R)$. The group $K^o(h)=K(h)\cap G^o(\R)$ is compact. Let $V$ be a representation of $G$ and let $\Psi$ be a $G$-invariant bilinear form on $V$. If $\Psi$ is a polarization of $(V,h_V)$ for one $h\in X$, then $\Psi$ is a polarization of $(V,h_V)$ for every $h\in X$. Indeed, putting $\Phi_h(x,y)=\Psi(x,C_hy)$, one has
\[
\tag{5.11.1}
  \Phi_{gh}=g(\Phi_h).
\]
We shall then say that $\Psi$ is a \underline{polarization of} $V$. By hypothesis, every representation of $G$ is polarizable.

\paragraph{5.12.}
Let $\Psi$ be a polarization of the adjoint representation of $G$ on $\Lie(G)$. By 5.10.2, $\Psi$ induces, at each point $h\in X$, a positive definite Hermitian form on
\[
\tag{5.12.1}
  (T_{\check X})_h
  \simeq
  F_h^{-\infty}/F_h^0(\Lie(G)_\C)
  \simeq
  \bigoplus_{p<0}\Lie(G)^{p,-p}.
\]
The polarization $\Psi$ therefore defines a \underline{Hermitian structure}, invariant under $G(\R)$, on $X$.

If $Z$ is the center of $G$, one already has
\[
\tag{5.12.2}
  F^{-\infty}/F^0\bigl(\Lie(G/Z)(\check X,\cO)\bigr)
  \xrightarrow{\sim} T_{\check X},
\]
and a polarization $\Psi$ of $\Lie(G/Z)$ already defines a structure of Hermitian variety on $X$. One may take for $\Psi$ the negative of the Killing form of $G/Z$.

Put
\[
\tag{5.12.3}
\begin{cases}
 (T_{\check X}^v)_h=\displaystyle\bigoplus_{p<0}(\Lie(G/Z))^{2p,-2p},\\[0.4em]
 (T_{\check X}^h)_h=\displaystyle\bigoplus_{p<0}(\Lie(G/Z))^{2p-1,-2p+1}.
\end{cases}
\]
Thus one obtains an orthogonal $C^\infty$ decomposition of the tangent bundle into two complex subbundles. One has
\[
\tag{5.12.4}
  (T_{\check X}^1)_h\subset (T_{\check X}^h)_h \qquad (h\in X).
\]

\paragraph{5.13.}
The subgroup $K(h)_\C\cap P(h)$ of $K(h)_\C$ is parabolic, so that $K(h)_\C\subset K(h)\cdot P(h)$ and $K(h)\cdot h=K(h)_\C\cdot h$ is a compact complex analytic subvariety of $X$. These varieties are the fibers of the $C^\infty$, $G(\R)$-equivariant fibration
\[
  X=G/H(h)\longrightarrow G/K(h).
\]
It is easily verified that $T_X^v$ is the tangent bundle to this fibration.

\paragraph{5.14.}
Recall that on a holomorphic Hermitian bundle $F$ there exists a unique Hermitian connection $\nabla$ such that $\nabla''=\partial''$. The \underline{curvature} $R$ of the bundle is, by definition, the curvature of this connection. It is a form of type $(1,1)$ with values in the Lie algebra of the unitary group of the bundle.

If $F$ is the constant bundle defined by a vector space $F_o$, and if the Hermitian form is written
\[
\tag{5.14.1}
  \Phi(x,y)=\Phi_o(hx,y),
\]
then, for $x:S\to F_o$ a $C^\infty$ section of $F$, one has
\[
\tag{5.14.2}
  \nabla x=\partial x+h^{-1}\partial' h\cdot x.
\]
It follows that, for $u$ and $v$ two tangent fields,
\[
\tag{5.14.3}
  R(u,v)=
  \partial_u(h^{-1}\partial'_v h)
  -\partial_v(h^{-1}\partial'_u h)
  -h^{-1}\partial'_{[uv]}h
  +[h^{-1}\partial'_u h,h^{-1}\partial'_v h].
\]

\paragraph{5.15.}
Let $F$ be the tangent bundle of a Hermitian complex analytic variety. For $u$ a tangent vector to $S$, denote by $u_{1,0}$ and $u_{0,1}=\overline{u_{1,0}}$ its components of type $(1,0)$ and $(0,1)$ in the complexification of the tangent bundle. The \underline{holomorphic curvature} of $S$ in the direction $u$ is the real number
\[
\tag{5.15.1}
  R(u)=
  \Phi(R(u_{1,0},u_{0,1})(u),u)/\Phi(u,u)^2.
\]

If $X$ is endowed with the Hermitian structure 5.12, one has the fundamental result ([8], §9):

\begin{theorem}[5.16]
\underline{There exists} $\lambda<0$ \underline{such that} $R(u)\le\lambda$ \underline{in every direction belonging to} $T_X^h$.
\end{theorem}

\underline{Proof} (a sketch). By homogeneity it is enough to prove 5.16 at a point $e\in X$. Put $H=H(e)$, $P=P(e)$, let $N^+$ be the unipotent radical of $P$, and let $N^-$ be the conjugate complex subgroup whose Lie algebra is
\[
  \mathfrak n^-=\bigoplus_{p<0}\Lie(G)^{p,-p}.
\]
In the calculations which follow, near $e$, we identify $X$ with $N^-$ by the maps
\[
  X=G/H \hookrightarrow G(\C)/P \hookleftarrow N^-.
\]
\underline{Left} translations on $N^-$ permit one to identify the tangent bundle with $\mathfrak n^-$. We shall denote by $t^+$, $t^o$, and $t^-$ the orthogonal projections of $\Lie(G)_\C$ onto $\mathfrak n^-$, $\Lie(H_\C)$, and $\mathfrak n^+=\Lie(N^+)$. For $g=np$ ($g\in G(\R)$, $n\in N^-$, $p\in P$), the differential, at $e$, of the action of $g$ on $X$ is identified with
\[
  t^-\ad p:\mathfrak n^-\longrightarrow \mathfrak n^- .
\]
The Hermitian structure is therefore written
\[
  \Phi_n(x,y)=\Phi_e(t^-\ad p^{-1}x,t^-\ad p^{-1}y).
\]
Denoting by $\Psi$ the polarization form and by $C$ the Cartan involution associated to $e$, one has
\[
\begin{aligned}
 \Phi_n(x,y)
 &=\Psi(t^-\ad p^{-1}x,\,t^- C\,\overline{\ad p^{-1}y})\\
 &=\Psi(t^-\ad p^{-1}x,\,\ad C(\overline p)^{-1}\cdot C\overline y)\\
 &=\Psi(\ad C(\overline p)\cdot t^-\cdot\ad p^{-1}\cdot x,\,C\overline y)\\
 &=\Phi_e(\ad C(\overline p)\cdot t^-\cdot\ad p^{-1}\cdot x,\,y).
\end{aligned}
\]
One verifies that, for $n=\exp(u)$ with $u\sim0$, \underline{to third order} one has
\[
  p=\exp\bigl(\overline u-\tfrac12t^o[u,\overline u]-t^+[u,\overline u]\bigr)
\]
and
\[
  h=\ad C(\overline p)\,t^-\,\ad p^{-1}=\exp(H),
\]
for
\[
  H=t^-\circ\ad(Cu-\overline u+[u,\overline u])
  -\tfrac12[t^-\ad Cu,t^-\ad\overline u].
\]

Under the hypotheses of 5.14, if $h=\exp(H)$ and if $H=0$ at $s_o\in S$, one deduces from 5.14.3 that, at $s_o$,
\[
\begin{aligned}
R(u,v)
={}&(\partial_u\partial'_v-\partial_v\partial'_u)H
-\tfrac12[\partial''_uH,\partial'_vH]
+\tfrac12[\partial''_vH,\partial'_uH]\\
&+(\partial'_u\partial'_v-\partial'_v\partial'_u-\partial'_{[uv]})H .
\end{aligned}
\]
Applying this formula, one sees that the curvature is given, at $e$, by
\[
  R(u,v)_e=S(u,\overline v)-S(v,\overline u),
\]
with
\[
\begin{aligned}
S(u,\overline v)
&=-t^-\circ\ad[u,\overline v]
+\tfrac12[t^-\ad Cu,t^-\ad\overline v]
+\tfrac12[-t^-\ad\overline v,t^-\ad Cu]\\
&=[t^-\ad Cu,t^-\ad\overline v]
-t^-\circ\ad[u,\overline v].
\end{aligned}
\]
For any $u,v\in\mathfrak n^-$, if $u_{10}$ and $u_{01}$ are the components of type $(10)$ and $(01)$ of $u$, as in 5.15, then, by the invariance of $\Psi$,


\[
\begin{aligned}
\Phi(R(u_{10},u_{01})(v),v)
&=\Psi(S(u,\overline u)(v),C\overline v)\\
&=\Psi\bigl([Cu,t^-[\overline u,v]]-t^-[\overline u[Cu,v]]-t^-[[u\overline u],v],C\overline v\bigr)\\
&=-\Psi\bigl(t^-[\overline u,v],[Cu,C\overline v]\bigr)
  +\Psi\bigl([Cu,v],[\overline u,C\overline v]\bigr)
  -\Psi\bigl([u,\overline u],[v,C\overline v]\bigr).
\end{aligned}
\]
For $u=v\in T_X^h$, one has $Cu=-u$ and
\[
\begin{aligned}
\Phi(R(u_{10},u_{01})(u),u)
&=\Psi\bigl(t^-[u,\overline u],[u,\overline u]\bigr)
  +\Psi\bigl([u,\overline u],[v,\overline v]\bigr)\\
&=-\Phi\bigl(t^-[u,\overline u],t^-[u,\overline u]\bigr)
  -\Phi\bigl([u,\overline u],[u,\overline u]\bigr)<0,
\end{aligned}
\]
which proves 5.16.

A generalization of Schwarz's lemma allows one to deduce from 5.16 (see [1], III, no. 10):

\begin{corollary}[5.17]
\underline{Let} $f:D\to X$ \underline{be a holomorphic map from the unit disk into $X$}. \underline{Endow} $D$ \underline{with the hyperbolic metric} $d$ \underline{for which it has curvature} $-1$, \underline{and} $X$ \underline{with the metric 4.12}. \underline{Then}, \underline{if} $f(D)$ \underline{is tangent to} $T_X^h$, \underline{one has}
\[
  d(f(x),f(y))\le |\lambda|\,d(x,y).
\]
\end{corollary}

\section{\underline{The Monodromy Theorem, after A. Borel}.}

\paragraph{6.1.}
Let $S$ be a connected complex analytic variety, endowed with a base point $s_o\in S$, and let $p:\widetilde S\to S$ be its universal covering. Let $H$ be a polarized family of Hodge structures of weight $n$ on $S$. Denote by $H_o$ the polarized Hodge structure which is the fiber of $H$ at $s_o$.

The polarization $\Psi$ is a bilinear form on $H_{o\Q}$. Denote by $G^o$ the algebraic group, defined over $\Q$, which is the neutral component of the group of automorphisms of $(H_{o\Q},\Psi)$; by $\varepsilon\in G^o$ the homothety of ratio $(-1)^n$; by $G$ the group deduced, by the procedure 5.5, from $(G^o,\varepsilon)$; and by $\Gamma$ the discrete group of automorphisms of $(H_{o\Z},\Psi)$. One has
\[
\tag{6.1.1}
  T:\pi_1(S,s_o)\longrightarrow\Gamma .
\]

\paragraph{6.2.}
For each point $s\in\widetilde S$, one has a natural isomorphism
\[
  H_{o\Z}\simeq H_{p(s),\Z},
\]
and the Hodge structure $H_{p(s)}$ is identified with a Hodge structure on $H_{o,\Z}$, defined by $h(s):\bS\to G_\R$. The morphisms $h(s)$ are positive and mutually conjugate, because they form a connected continuous family and $G$ is reductive. If $X$ is the corresponding moduli space, defined in § 5, one therefore has a map
\[
  h:\widetilde S\longrightarrow X.
\]
The family $H$ is the pullback along $h$ of the system of Hodge structures on $X$ defined in (5.9) by the natural weight-$n$ representation of $G$ on $H_o$.

The morphism $h$ is equivariant, via (5.5.1), for the actions of $\pi_1(S,s_o)$ and $\Gamma$ on $\widetilde S$ and $X$.

By 3.4 and 3.6, the map $h$ is holomorphic and $\operatorname{Im}(dh)\subset T_X^1$ (cf. 5.10).

\paragraph{6.3.}
Suppose that $S$ is the punctured unit disk
\[
  D^*=\{z\mid 0<|z|<1\}
\]
and take for the universal covering of $D^*$ the Poincaré half-plane $Y$, endowed with
\[
  \exp(2\pi i y):Y\longrightarrow D^*.
\]
If $T\in\Gamma$ is the monodromy transformation (4.3) of $H_\Z$, the map $h$ satisfies
\[
\tag{6.3.1}
  h(y+1)=T h(y).
\]
Choose $h_o\in X$. If $h(y)=g_y\cdot h_o$, the distance in $X$, for a metric (5.12), satisfies
\[
\tag{6.3.2}
  d(h(y+1),h(y))=d(Tg_yh_o,g_yh_o)=d(g_y^{-1}Tg_yh_o,h_o).
\]
By 5.17, if $Y$ is endowed with its natural hyperbolic distance, the map $h$ decreases distances (suitably normalized). In $Y$, one has
\[
  d(y,y+1)\sim (\operatorname{Im}(y))^{-1}
\]
and therefore
\[
\tag{6.3.3}
  d(g_y^{-1}Tg_yh_o,h_o)\longrightarrow 0\qquad\text{for }\operatorname{Im}(y)\to\infty .
\]
The same argument applies to $T^2$; moreover, $T^2\in G(\R)$, and (6.2.3) implies that a limit of conjugates of $T^2$ lies in the compact subgroup $H(h_o)$ of $G(\R)$. The eigenvalues of $T^2$ (and of $T$) are therefore all of absolute value $1$.
These eigenvalues form a set of algebraic integers stable under conjugation. Now,

\begin{lemma}[6.4]
\underline{If all the complex conjugates of an algebraic integer} $\alpha$ \underline{have absolute value} $1$, \underline{then} $\alpha$ \underline{is a root of unity}.
\end{lemma}

The eigenvalues of $T$ are therefore roots of unity, which proves

\begin{monodromyTheorem}[6.5 (A. Borel)]
\underline{Let} $H$ \underline{be a polarized family of Hodge structures on the punctured disk} $D^*$. \underline{There exist integers} $N$ \underline{and} $M$ \underline{such that the monodromy transformation} $T$ \underline{satisfies}
\[
  (T^N-I)^M=0.
\]
\end{monodromyTheorem}

\begin{center}
\bfseries BIBLIOGRAPHY
\end{center}

\begin{enumerate}[label={[\arabic*]},leftmargin=2.8em]
\item P. A. GRIFFITHS - \uline{Periods of integrals on algebraic manifolds}.\\
I (Construction and properties of the modular Varieties), Am. J. Math., XC 2, 1968, p. 568-626.\\
II, Am. J. Math., XC 3, 1968, p. 805-865.\\
III (Some global differential-geometric properties of the period mapping), to appear in Publ. I.H.E.S.

\item P. A. GRIFFITHS - \uline{Some results on Moduli and Periods of Integrals on Algebraic Manifolds III}, mimeographed notes from Princeton.

\item P. A. GRIFFITHS - \uline{On the periods of integrals on algebraic manifolds}, Rice un. studies, \uline{54}, 4, 1968.\\
This article summarizes, without proof, [1] I, II.

\item P. A. GRIFFITHS - \uline{On the periods of certain rational integrals}, I, II, Annals of Maths., 90 (1969), p. 460-541. III, to appear in Annals of Maths.

\item P. A. GRIFFITHS - \uline{A theorem on periods of integrals on algebraic manifolds}, mimeographed notes from Yale.

\item P. A. GRIFFITHS - \uline{Some results on algebraic cycles on algebraic manifolds}, in Bombay Colloquium 1968, Oxford University Press.

\item P. A. GRIFFITHS - \uline{Periods of integrals on algebraic manifolds} (Summary of main results and discussions of open problems and conjectures), Bull. Am. Math. Soc., \uline{75}, 2, (1970), p. 228-296.\\
This article is very interesting because of the many conjectures it discusses.

\item P. A. GRIFFITHS and W. SCHMID - \uline{Locally homogeneous complex manifolds}, Acta Mathematica, 1970, p. 253-302.\\
This very readable article contains much information on the spaces defined in § 5 and the cohomology of their quotients by discrete subgroups.

\item N. KATZ - \uline{Nilpotent connections and the monodromy theorem, Applications of a results of Turrittin}, to appear in Publ. I.H.E.S.

\item N. KATZ and T. ODA - \uline{On the differentiation of De Rham cohomology classes with respect to parameters}, J. Math. Kyoto Univ., \uline{8}, 2, 1968, p. 199-213.
\end{enumerate}

\end{document}
